\chapter{Arquitectura del sistema}
\section{Diseño general del sistema}

Komuness sigue una arquitectura cliente-servidor de tipo \textbf{N-capas}, con separación explícita entre presentación, lógica de negocio y persistencia. El backend se implementa en \textbf{Node.js} con el framework \textbf{Express} \cite{expressjs2024} y \textbf{TypeScript}, organizado en las capas \texttt{routes} (define endpoints y aplica middlewares de autenticación/autorización), \texttt{controllers} (lógica de negocio y validación de cada operación), \texttt{models} (esquemas de \textbf{Mongoose} \cite{mongoosejs2024} sobre \textbf{MongoDB}), \texttt{interfaces} (tipos de TypeScript compartidos), \texttt{middlewares} (autenticación JWT, verificación de roles, límites de publicación y carga de archivos con Multer) y \texttt{utils} (conexión a base de datos, envío de correo, JWT, almacenamiento en disco y GridFS, y verificación de archivos). El frontend se implementa como una \textit{Single Page Application} en \textbf{React} \cite{reactjs2024}, con enrutamiento de cliente (\texttt{react-router-dom}), manejo de estado de sesión mediante \textit{React Context} (\texttt{AuthContext}, \texttt{ThemeContext}) y consumo de la API a través de \texttt{fetch} contra la URL configurada en \texttt{REACT\_APP\_BACKEND\_URL}.

La autenticación se basa en tokens JWT \cite{jwt2015} enviados en el encabezado \texttt{Authorization: Bearer}. El middleware \texttt{authMiddleware} valida el token, recupera el usuario real desde la base de datos (nunca confía únicamente en el contenido del token) y aplica automáticamente el \textit{downgrade} de un usuario premium cuya membresía haya vencido. Sobre esa base, el middleware \texttt{verificarRoles} restringe el acceso a operaciones administrativas según el campo \texttt{tipoUsuario} (0 = super-administrador, 1 = administrador, 2 = básico, 3 = premium). El almacenamiento de archivos combina dos estrategias: el \textbf{disco del servidor} (directorio \texttt{UPLOAD\_DIR}, por defecto \texttt{/srv/uploads}) para la biblioteca documental, los archivos de perfil y las imágenes de ``Acerca de'' y la temática, y \textbf{GridFS} de MongoDB para las imágenes de publicaciones y publicidad, servidas por \texttt{/api/files/:id}. El módulo \texttt{utils/digitalOceanSpace.ts} conserva ese nombre por razones históricas, pero hoy solo escribe en disco; las dependencias \texttt{aws-sdk} y \texttt{multer-s3} siguen declaradas en \texttt{package.json} aunque el código ya no las usa. La biblioteca no se publica como directorio estático: sus archivos se descargan por \texttt{GET /api/biblioteca/files/:id}, de modo que los archivos pendientes o rechazados no quedan expuestos. El envío de correos (confirmación de cuenta, recuperación de contraseña, boletines) se realiza con \textbf{Nodemailer}, y los pagos de membresías premium se procesan con el SDK de \textbf{PayPal} \cite{paypaldev2024}, incorporando una capa de reintentos con \textit{backoff} exponencial ante errores transitorios del proveedor.

\subsection{Diseño de la persistencia}
La persistencia se modela con \textbf{Mongoose} sobre \textbf{MongoDB}, por lo que el modelo de datos es orientado a documentos: la mayoría de las relaciones se expresan como referencias (\texttt{ObjectId} con \texttt{ref}) resueltas mediante \texttt{populate}, y solo los datos que siempre se leen junto a su documento padre (comentarios de una publicación, formación académica de un perfil, miembros del equipo en la sección ``Acerca de'') se embeben como subdocumentos. La Figura \ref{ERNucleo} muestra el modelo conceptual de las entidades centrales del dominio (usuarios, contenido y actividades), y la Figura \ref{ERSoporte} muestra las entidades de soporte administrativo y comercial.

\begin{figure}[H]
  \centering
  \resizebox{\textwidth}{!}{%
  \begin{tikzpicture}[
    box/.style={draw, rectangle, rounded corners=2pt, minimum width=2.6cm, minimum height=0.9cm, align=center, font=\small, inner sep=4pt},
    every path/.style={-{Latex[length=2mm]}},
    lbl/.style={font=\scriptsize, fill=white, inner sep=1pt}
  ]
    \node[box] (usuario) at (0,8) {Usuario};
    \node[box] (perfil) at (0,4) {Perfil};
    \node[box] (catpref) at (5.5,8) {CategoriaPreferencia};
    \node[box] (categoria) at (5.5,12.5) {Categoria};
    \node[box] (publicacion) at (12,8) {Publicacion \\ \scriptsize(+Comentario embebido)};
    \node[box] (actividad) at (12,4) {Actividad};
    \node[box] (notificacion) at (5.5,4) {Notificacion \\ \scriptsize(+eliminadasPor)};
    \node[box] (cuenta) at (17.5,8) {CuentaRegresiva};

    \draw[-] (usuario) -- (perfil) node[lbl, midway]{1:1};
    \draw[-] (usuario) -- (catpref) node[lbl, midway]{1:N};
    \draw[-] (categoria) -- (catpref) node[lbl, midway]{1:N};
    \draw[-] (categoria) -- (publicacion) node[lbl, midway]{1:N};
    \draw[-] (usuario) to[bend left=15] node[lbl, midway]{1:N (autor)} (publicacion);
    \draw[-] (usuario) to[bend right=15] node[lbl, midway]{1:N} (actividad);
    \draw[-] (usuario) -- (notificacion) node[lbl, midway]{N:M (recipientes)};
    \draw[-, dashed] (publicacion) -- (notificacion) node[lbl, midway]{1:N (opcional)};
    \draw[-, dashed] (actividad) to[bend left=10] node[lbl, midway]{genera} (notificacion);
    \draw[-, dashed] (publicacion) to[loop above] node[lbl, above]{origen (publicación automática)} (publicacion);
    \draw[-, dashed] (publicacion) -- (cuenta) node[lbl, midway]{0..1 (evento)};
  \end{tikzpicture}}
  \caption{\textbf{Modelo conceptual de entidades centrales: usuarios, contenido, actividades y cuentas regresivas}}
  \label{ERNucleo}
\end{figure}

\begin{figure}[H]
  \centering
  \resizebox{\textwidth}{!}{%
  \begin{tikzpicture}[
    box/.style={draw, rectangle, rounded corners=2pt, minimum width=2.6cm, minimum height=0.9cm, align=center, font=\small, inner sep=4pt},
    every path/.style={-{Latex[length=2mm]}},
    lbl/.style={font=\scriptsize, fill=white, inner sep=1pt}
  ]
    \node[box] (usuario) at (0,11) {Usuario};
    \node[box] (publicacion) at (0,1) {Publicacion};
    \node[box] (payment) at (6,13) {Payment};
    \node[box] (boletin) at (6,8.5) {Boletin};
    \node[box] (publicidad) at (6,4) {Publicidad};
    \node[box] (archivo) at (12,11) {Archivo};
    \node[box] (folder) at (12,6.5) {Folder};
    \node[box] (categoriaarch) at (16.5,11) {Categoria};
    \node[box] (catalogo) at (12,1) {\begin{tabular}{c} Catálogo institucional \\ \scriptsize Tutorial, ProyectoDestacado, \\ \scriptsize Etiqueta, Configuracion, \\ \scriptsize SeccionAcerca, PaqueteSuscripcion \end{tabular}};

    \draw[-] (usuario) -- (payment) node[lbl, midway]{1:N};
    \draw[-] (usuario) -- (boletin) node[lbl, midway]{1:N (creadoPor)};
    \draw[-, dashed] (usuario) to[bend right=25] node[lbl, midway]{N:M (destinatarios)} (boletin);
    \draw[-] (usuario) -- (publicidad) node[lbl, midway]{1:N (autor)};
    \draw[-, dashed] (publicacion) -- (publicidad) node[lbl, midway]{1:N (opcional)};
    \draw[-] (usuario) -- (archivo) node[lbl, midway]{1:N (autor)};
    \draw[-] (folder) -- (archivo) node[lbl, midway]{1:N};
    \draw[-] (folder) to[loop left] node[lbl, left]{1:N (padre)} (folder);
    \draw[-] (categoriaarch) -- (archivo) node[lbl, midway]{1:N};
  \end{tikzpicture}}
  \caption{\textbf{Modelo conceptual de soporte administrativo y comercial}}
  \label{ERSoporte}
\end{figure}

Las entidades del catálogo institucional (\texttt{Tutorial}, \texttt{ProyectoDestacado}, \texttt{Etiqueta}, \texttt{Configuracion}, \texttt{SeccionAcerca} y \texttt{PaqueteSuscripcion}) son mayormente independientes entre sí: se consultan como catálogos administrados desde el panel de administración y no participan de relaciones de cardinalidad múltiple con el resto del dominio, salvo referencias puntuales de auditoría (por ejemplo, \texttt{Configuracion.actualizadoPor} apunta a \texttt{Usuario}).

\subsection{Diagrama de clases}
La Figura \ref{DiagramaClases} presenta el diagrama de clases simplificado del núcleo del dominio, con los atributos más relevantes de cada esquema de Mongoose. Se destaca que \texttt{Usuario} es el eje central del modelo: de él dependen el \texttt{Perfil} (relación uno a uno), las \texttt{Publicacion} que autora (uno a muchos) y las \texttt{Actividad} que crea (uno a muchos). En el Sprint 2, \texttt{Publicacion} incorporó una referencia opcional a sí misma (\texttt{publicacionOrigenId}), que enlaza la publicación generada automáticamente con el evento o emprendimiento del que proviene, y \texttt{Perfil} incorpora \texttt{provincia}, \texttt{canton} y \texttt{perfilPublico}, usados por los filtros y la elegibilidad del banco de profesionales. En el Sprint 3 el mismo campo \texttt{publicacionOrigenId} pasa a apuntar también al tutorial, proyecto destacado o archivo de biblioteca de origen (que son otras colecciones, por lo que la referencia \texttt{ref: 'Publicacion'} es solo informativa), y \texttt{categoria} deja de ser obligatoria cuando existe \texttt{publicacionOrigenId}. Además: \texttt{Usuario} incorpora \texttt{puedeCrearCuentaRegresiva}; \texttt{Notificacion} incorpora \texttt{eliminadasPor} (usuarios que la ocultaron para sí mismos); \texttt{Archivo} incorpora \texttt{categoria}; \texttt{Publicidad} incorpora \texttt{titulo}, \texttt{fechaInicioEvento} y \texttt{fechaFinEvento}; y se crea la entidad \texttt{CuentaRegresiva} (\texttt{titulo}, \texttt{fechaObjetivo}, \texttt{publicacion} única y opcional, \texttt{enlace}, \texttt{creadoPor}, \texttt{activo}). \texttt{Categoria} clasifica tanto las publicaciones como las preferencias de los usuarios.

\begin{figure}[H]
  \centering
  \resizebox{\textwidth}{!}{%
  \begin{tikzpicture}[
    every path/.style={-{Latex[length=2mm]}},
    lbl/.style={font=\scriptsize, fill=white, inner sep=1pt},
    classbox/.style={
      rectangle split, rectangle split parts=2, draw, rounded corners=1pt,
      rectangle split part align={left, left},
      align=left, font=\scriptsize, text width=3.6cm
    }
  ]
    \node[classbox] (usuario) at (0,7) {
      \textbf{Usuario}
      \nodepart{two}
      nombre, apellido\\ email, password\\ tipoUsuario: 0..3\\ fechaVencimientoPremium\\ emailConfirmado\\ puedeCrearCuentaRegresiva
    };
    \node[classbox] (perfil) at (0,1) {
      \textbf{Perfil}
      \nodepart{two}
      usuarioId\\ ocupacionPrincipal\\ provincia, canton\\ formacionAcademica[]\\ habilidades[]\\ enBancoProfesionales\\ perfilPublico
    };
    \node[classbox] (categoria) at (7,10) {
      \textbf{Categoria}
      \nodepart{two}
      nombre\\ estado
    };
    \node[classbox] (publicacion) at (7,7) {
      \textbf{Publicacion}
      \nodepart{two}
      titulo, contenido\\ autor, categoria\\ tag: publicacion\textbar evento\textbar emprendimiento\\ fecha, fechaEvento\\ precio, comunidad\\ ubicacion (lat, lng, direccion)\\ adjunto[], comentarios[]\\ publicacionOrigenId
    };
    \node[classbox] (cuenta) at (14,7) {
      \textbf{CuentaRegresiva}
      \nodepart{two}
      titulo, enlace\\ fechaObjetivo\\ publicacion (0..1, única)\\ creadoPor, activo
    };
    \node[classbox] (actividad) at (7,1) {
      \textbf{Actividad}
      \nodepart{two}
      nombre, descripcion\\ fecha, hora\\ ubicacionEscrita\\ ubicacion: GeoJSON Point\\ usuario
    };

    \draw (usuario) -- (perfil) node[lbl, midway]{1..1};
    \draw (usuario) -- (publicacion) node[lbl, midway, sloped]{1..*} node[lbl, pos=0.85]{autor};
    \draw (usuario) -- (actividad) node[lbl, midway, sloped]{1..*};
    \draw (categoria) -- (publicacion) node[lbl, midway]{1..*};
    \draw[dashed] (cuenta) -- (publicacion) node[lbl, midway]{0..1};
    \draw[dashed] (publicacion) to[loop right] node[lbl, right]{origen 0..1} (publicacion);
  \end{tikzpicture}}
  \caption{\textbf{Diagrama de clases simplificado del núcleo del dominio}}
  \label{DiagramaClases}
\end{figure}

\section{Diseño del frontend}
El frontend es una \textit{Single Page Application} en React (\texttt{FrontEnd/src}). El componente \texttt{Rutas} declara el enrutamiento de cliente y envuelve todas las pantallas con la barra de navegación (\texttt{Navbar}), el modal de encuesta inicial y el pie de página. La Tabla \ref{tab:rutas} resume las rutas principales.

\begin{table}[H]
\centering
\small
\begin{tabular}{|p{4.6cm}|p{10.2cm}|}
\hline
\textbf{Ruta} & \textbf{Pantalla / comportamiento} \\
\hline
\texttt{/} & Pantalla de inicio (\texttt{PantallaInicio}): carrusel de la portada (\texttt{CarruselPortada}), cuentas regresivas (\texttt{CuentasRegresivas}), publicidad (\texttt{PublicidadBanner}), contenido principal (\texttt{InicioPrincipal}) y sección ``Quiénes somos'' (\texttt{QuienesSomosSection}, ancla \texttt{\#quienes-somos}), que los administradores editan desde la misma página. \\
\hline
\texttt{/eventos}, \texttt{/publicaciones}, \texttt{/emprendimientos} & \texttt{LandingPage} con la etiqueta correspondiente; listado, filtros y, en eventos, el calendario. \\
\hline
\texttt{/eventos/:id}, \texttt{/publicaciones/:id}, \texttt{/emprendimientos/:id} & Detalle (\texttt{PublicacionDetalle}); la ruta conserva la sección de origen. \\
\hline
\texttt{/calendario}, \texttt{/acerca-de}, \texttt{/admin/acerca-de} & Redirigen a \texttt{/eventos}, a \texttt{/\#quienes-somos} y a \texttt{/}. \\
\hline
\texttt{/profesionales}, \texttt{/perfil/:id}, \texttt{/mi-perfil/editar} & Banco de profesionales, perfil público y edición del perfil. \\
\hline
\texttt{/biblioteca/:id}, \texttt{/notificaciones}, \texttt{/checkout-premium} & Biblioteca por carpeta, notificaciones en vista móvil y pago de membresía. \\
\hline
\texttt{/admin/categorias}, \texttt{/admin/boletines} & Pantallas de administración. \\
\hline
\texttt{/cuentas-regresivas} & Gestión de cuentas regresivas (\texttt{AdminCuentasRegresivas}); visible para administradores y usuarios con el permiso. \\
\hline
\end{tabular}
\caption{\textbf{Rutas principales del frontend}}
\label{tab:rutas}
\end{table}

Las URL del backend se resuelven en \texttt{src/utils/api.js}: en desarrollo se usa \texttt{REACT\_APP\_BACKEND\_URL} (y no un puerto fijo, para evitar el conflicto del puerto 5000 con AirPlay en macOS). El módulo \texttt{src/utils/secciones.js} centraliza la correspondencia etiqueta--sección (\texttt{evento} $\rightarrow$ \texttt{/eventos}, \texttt{emprendimiento} $\rightarrow$ \texttt{/emprendimientos}, \texttt{publicacion} $\rightarrow$ \texttt{/publicaciones}) mediante \texttt{getRutaDetalle}, que usan las tarjetas, los carruseles, el modal de evento y el buscador. Los filtros de las listas (categoría, fechas, precio, comunidad y búsqueda) viven en los parámetros de la URL, de modo que el botón ``Borrar filtros'' solo debe navegar a la ruta sin parámetros y los controles se resincronizan.

\section{Diseño de los módulos del Sprint 2}

\subsection{Publicación automática al crear un evento o emprendimiento}
La creación con adjuntos (\texttt{POST /api/publicaciones/v2}, función \texttt{createPublicacionA}) guarda primero el documento solicitado. Si su etiqueta es \texttt{evento} o \texttt{emprendimiento}, genera además una segunda publicación (etiqueta \texttt{publicacion}) con el mismo título, contenido, categoría, comunidad, fecha y estado de publicación. Las decisiones de diseño relevantes son:
\begin{itemize}
    \item \textbf{Copia independiente de adjuntos:} cada archivo se vuelve a guardar en GridFS con una clave nueva, porque la eliminación de una publicación borra sus adjuntos y, de compartirlos, rompería las imágenes de la otra.
    \item \textbf{Sincronización (Sprint 3):} en el Sprint 2 la relación era solo informativa. Desde el Sprint 3 el original y su copia se editan, aprueban y eliminan juntos (ver la sección de publicaciones enlazadas más adelante); los adjuntos siguen siendo independientes.
    \item \textbf{Tolerancia a fallos:} el paso se ejecuta dentro de un \texttt{try/catch}; si falla, solo se registra una advertencia y el evento o emprendimiento queda creado.
\end{itemize}

\subsection{Ubicación en mapa}
La validación de la ubicación se extrajo a \texttt{backend/src/utils/ubicacion.ts} (\texttt{validarUbicacion}) y la comparten la creación, la edición administrativa y la solicitud de edición. Acepta un objeto o un JSON en texto (por \textit{FormData}), valida rangos de latitud y longitud y la dirección, y genera el enlace \texttt{mapLink} a OpenStreetMap. En el cliente, \texttt{MapaUbicacion} (Leaflet) permite elegir el punto y obtener la dirección con el servicio Nominatim; se muestra tanto para eventos como para emprendimientos. Nótese que las publicaciones guardan la ubicación como \texttt{latitude}/\texttt{longitude}, mientras que \texttt{Actividad} usa un punto GeoJSON con índice \texttt{2dsphere}.

\subsection{Elegibilidad del banco de profesionales}
La regla de qué campos debe tener un perfil para unirse al banco se centralizó en \texttt{backend/src/utils/perfilElegibilidad.ts} (\texttt{obtenerFaltantesBanco}). Se usa en tres lugares: el endpoint de consulta de elegibilidad, el alternado de membresía en el banco y el guardado del perfil cuando se solicita \texttt{enBancoProfesionales}. De esta forma el servidor no depende de la validación del frontend. En los dos últimos casos, un perfil incompleto produce una respuesta con \texttt{success: false} y la lista de campos faltantes en \texttt{data}.

\subsection{Búsqueda de profesionales}
\texttt{obtenerProfesionales} construye los patrones de búsqueda con \texttt{patronInsensibleAcentos}, que reemplaza cada letra acentuable por una clase de caracteres (por ejemplo, \texttt{[eéèëêEÉÈËÊ]}) y escapa el resto de los caracteres, evitando expresiones regulares inválidas. La paginación (\texttt{page}, \texttt{limit}) la consume el frontend con el botón ``Cargar más'', que acumula las páginas recibidas.

\section{Diseño de los módulos del Sprint 3}

\subsection{Servicio de publicación automática y publicaciones enlazadas}
La creación de la publicación automática se centralizó en \texttt{services/publicacionAutomatica.service.ts} (\texttt{crearPublicacionAutomatica}), que recibe el título, la descripción, la URL, el autor, el identificador del origen (\texttt{publicacionOrigenId}) y el tipo de origen (\texttt{evento}, \texttt{emprendimiento}, \texttt{tutorial}, \texttt{proyecto-destacado} o \texttt{biblioteca}). Genera una publicación con etiqueta \texttt{publicacion}, un enlace externo cuyo texto depende del origen (``Ver tutorial'', ``Ver proyecto'', ``Ver recurso''), el contenido breve recortado a 180 caracteres y, cuando ya es visible, una notificación general mediante \texttt{notificarPublicacionAutomatica}. La categoría es opcional solo para estas publicaciones: el esquema de \texttt{Publicacion} la exige únicamente si no hay \texttt{publicacionOrigenId}, y los controladores eliminan ese campo de lo que envía el cliente para que nadie pueda saltarse la categoría obligatoria. Eventos y emprendimientos (que sí tienen categoría) crean su copia dentro de \texttt{createPublicacionA}; tutoriales, proyectos destacados y archivos de biblioteca usan el servicio.

Para eventos y emprendimientos, \texttt{utils/publicacionEnlazada.ts} implementa la edición y eliminación enlazadas (BKD67). \texttt{CAMPOS\_ENLAZADOS} lista los campos compartidos; después de guardar un cambio, \texttt{sincronizarPublicacionesEnlazadas} copia al otro documento los campos modificados y recalcula su fecha de expiración, y \texttt{findPublicacionesEnlazadas} permite que \texttt{deletePublicacion} elimine ambos documentos y sus adjuntos. Como \texttt{publicado} es un campo compartido, aprobar uno aprueba el otro; la notificación general se envía una sola vez (BKD68) y, para evitar avisos duplicados, no se envía la notificación a los seguidores de la categoría cuando el origen es un evento o emprendimiento. La tarea de expiración (\texttt{purgeExpiredPublicaciones}) elimina también las copias de los eventos vencidos. Las publicaciones automáticas no cuentan contra el límite de publicaciones del usuario ni aparecen como solicitudes separadas en la lista de pendientes.

\subsection{Biblioteca: categorías y publicación automática}
\texttt{Archivo} incorpora el campo \texttt{categoria} (referencia a \texttt{Categoria}, obligatoria al subir pero opcional en el esquema para no invalidar archivos antiguos). \texttt{uploadFiles} valida la categoría, acepta nombres personalizados (\texttt{resolveFileName}, que conserva la extensión original) y guarda cada binario en disco con un nombre aleatorio, independiente del nombre visible. La función \texttt{publicarArchivoAutomatico} se invoca cuando un archivo queda aprobado, ya sea al subirlo un administrador o en \texttt{approveFile}; comprueba que no exista una publicación con ese origen y nunca lanza una excepción. \texttt{deleteFileById} elimina la publicación asociada. Los endpoints de listado y búsqueda reciben el parámetro \texttt{categoria} y la interfaz (\texttt{biblioteca.js}) ofrece el selector obligatorio al subir y un filtro de búsqueda.

\subsection{Cuentas regresivas}
Las cuentas regresivas de la página principal se componen de la entidad \texttt{CuentaRegresiva}, el servicio \texttt{services/cuentaRegresiva.service.ts}, las utilidades \texttt{utils/cuentaRegresiva.ts} y los controladores y rutas homónimos. Todas las fechas se interpretan en hora de Costa Rica (UTC$-6$, sin horario de verano): \texttt{construirFechaObjetivo} convierte la fecha y la hora de un evento en el instante real y rechaza fechas que el motor ``corregiría'' (por ejemplo, 30 de febrero). \texttt{obtenerActivas} devuelve como máximo \texttt{MAX\_CUENTAS\_ACTIVAS} $=4$ contadores. Un contador puede ser libre o seguir a un evento mediante una referencia única y opcional a la publicación; en ese caso \texttt{sincronizarCuentaRegresivaDeEvento} actualiza título y fecha cuando el evento cambia, lo elimina si el evento deja de serlo o pierde su fecha, y se invoca también al aprobar o editar la copia enlazada. La regla de permiso (\texttt{puedeGestionarCuentasRegresivas}) existe de forma idéntica en el backend y en el frontend (\texttt{utils/permisos.js}), pero solo la del backend es la que protege la API.

\subsection{Notificaciones por usuario}
En lugar de borrar el documento, \texttt{deleteNotificacion} agrega al usuario al arreglo \texttt{eliminadasPor} (\texttt{\$addToSet}) y \texttt{deleteAllNotificaciones} lo hace para todas las notificaciones visibles para él. Las consultas (\texttt{getNotificaciones}, marcar como vista) comparten \texttt{filtroVisibleParaUsuario}, excluyen las eliminadas por el usuario y no devuelven el campo \texttt{eliminadasPor}. Así una notificación general (\texttt{recipientes} vacío) sigue visible para el resto de la comunidad.

\subsection{Seguridad de las cargas de archivos y de la actualización de usuarios}
\texttt{utils/fileContent.ts} (\texttt{verifyFileContent}) lee los bytes de cada archivo, sin dependencias externas, para comprobar que su tipo real coincide con su extensión y que no contiene ejecutables, scripts ni contenido activo conocido; cada ruta que recibe archivos declara una política (\texttt{images}, \texttt{pdf} y \texttt{library}). \texttt{utils/uploadSecurity.ts} (\texttt{secureUploads}) se coloca justo después de multer: rechaza con 400 si algún archivo no es válido, reemplaza el MIME declarado por el real y borra del disco todos los archivos si la respuesta termina en error o la conexión se corta. El manejador global de errores traduce \texttt{ImageValidationError} a una respuesta 400 con código \texttt{INVALID\_FILE\_TYPE}. Por otra parte, \texttt{updateUsuario} aplica una lista blanca de campos (\texttt{CAMPOS\_EDITABLES\_USUARIO}) y la jerarquía entre roles, de modo que los cambios sensibles solo se hacen por las rutas dedicadas (\texttt{/membresia}, \texttt{/limite}, \texttt{/premium-vencimiento}, \texttt{/cuenta-regresiva}).

\subsection{Portada, publicidad y envío de correo}
El frontend reemplazó la pantalla separada ``Acerca de'' por la edición dentro de la página de inicio: \texttt{CarruselPortada} muestra y reordena las imágenes del carrusel (el reordenamiento se envía en \texttt{imagenesProyectos} y el servidor solo lo acepta si es una permutación exacta de las imágenes existentes) y \texttt{PublicidadBanner} presenta el título, la descripción y las fechas de la publicidad. El controlador de publicidad exige ahora que \texttt{PUBLIC\_BASE\_URL} esté definida y falla con un error claro en lugar de usar una dirección fija. \texttt{utils/mail.ts} usa la cuenta de Gmail con contraseña de aplicación cuando \texttt{MAIL\_USER} termina en \texttt{@gmail.com} y, en otro caso, el servidor SMTP definido por \texttt{MAIL\_HOST} y \texttt{MAIL\_PORT}.
